% ============================================================
% P18 — Voucher Pelanggan
% ============================================================
\flowmeta{P18}{Voucher Pelanggan}{Pelanggan}

\begin{flowsection}{Tujuan}
Melihat dan mengklaim voucher promo, serta memakainya saat checkout.
\end{flowsection}

\begin{flowsection}{Prasyarat}
Sudah login sebagai pelanggan.
\end{flowsection}

\begin{flowsection}{Langkah Mengelola Voucher}
\begin{enumerate}
  \item Buka \texttt{/profile} lalu tab \textbf{Voucher Diskon}.
  \item Lihat daftar voucher yang tersedia/aktif.
  \item Klik \textbf{Klaim} untuk menyimpan voucher ke akun.
  \item Voucher yang dimiliki dapat dimasukkan sebagai kode kupon di keranjang/checkout.
\end{enumerate}
\end{flowsection}

\shot{gambar/pelanggan/p18-1.png}{Tab Voucher Diskon pada halaman profil pelanggan}{fig:p18-1}

\begin{flowsection}{Hasil yang Diharapkan}
Voucher tersimpan pada akun dan dapat diterapkan pada transaksi yang memenuhi
syarat (minimal belanja/kategori).
\end{flowsection}

\begin{flowsection}{Penanganan Error dan Kasus Khusus}
\begin{itemize}
  \item Kuota voucher habis $\rightarrow$ klaim ditolak.
  \item Tidak memenuhi syarat $\rightarrow$ muncul pesan kesalahan kupon.
\end{itemize}
\end{flowsection}

\begin{flowsection}{Kaitan Modul}
FE \texttt{ProfilePage.jsx} (tab voucher), \texttt{authService.claimVoucher}. BE
\texttt{VoucherController}. Rute \texttt{GET /api/vouchers}, \texttt{POST /api/vouchers/claim},
\texttt{POST /api/vouchers/validate}.
\end{flowsection}

\docver{v1.0}{Sep 2026}
